\documentclass[11pt]{letter}
\usepackage[margin=1in]{geometry}
\usepackage{amsmath}              % \text{} inside math mode (subscripts)
\usepackage{hyperref}

% Sender information
\address{
Dr. Ram Chand\\
Chairperson, Department of Natural Sciences\\
The Begum Nusrat Bhutto Women University\\
Sukkur, Sindh, Pakistan\\
Email: ram.chand@bnbwu.edu.pk
}

\signature{Dr. Ram Chand}

\begin{document}

\begin{letter}{
Editor-in-Chief\\
Microfluidics and Nanofluidics
}

\opening{Dear Editor,}

I am pleased to submit my research manuscript for publication in
\textit{Microfluidics and Nanofluidics}, entitled
\textbf{``A deformability threshold for stiffness-based sorting in
deterministic-lateral-displacement microfluidic devices: a 2D
lattice-Boltzmann immersed-boundary parameter sweep''}.

Deterministic lateral displacement (DLD) is the workhorse passive
sorting technology of modern microfluidic cell-separation chips,
but the founding Huang--Holm critical-bin theory was derived for
rigid spheres. Experimental and computational studies have since
shown that deformable particles can be sorted by stiffness at
\emph{identical particle size}, opening DLD to label-free
mechanophenotyping of red blood cells, circulating tumour cells
and synthetic vesicle drug carriers. What was missing from the
literature is a quantitative answer to the design question every
DLD chip designer faces: \emph{at what Capillary number does
stiffness-based sorting actually switch on?} Below that threshold
the chip behaves geometrically (rigid-sphere DLD) and stiffness
contrast does nothing for the user; above it, sorting fidelity
grows monotonically with contrast. This manuscript identifies that
threshold quantitatively.

The main contributions of this work are:
\begin{itemize}
  \item \textbf{Custom lattice-Boltzmann immersed-boundary code with
        full discretisation documented:} all simulations were
        performed with \textsc{SoftFlow}, a C++ LBM-IBM code I
        developed in-house. The manuscript reports the complete
        numerical scheme --- D2Q9 distribution functions, regularised
        BGK collision (Latt--Chopard), Guo body forcing, interpolated
        bounce-back (Bouzidi--Firdaouss--Lallemand) for curved
        pillars, multi-direct-forcing IBM (Luo et al.), and the
        Skalak constitutive law for the membrane --- so that the
        results are reproducible from the equations alone.

  \item \textbf{Matched-membrane no-contrast diagonal control:} the
        two species in the sweep differ \emph{only} in their surface
        shear modulus $G_s$; bending, area and perimeter moduli are
        held identical. The diagonal of the sweep grid
        ($G_s^{\text{soft}} = G_s^{\text{stiff}}$) is therefore a
        strict no-contrast control in which the displacement-angle
        separation must vanish in expectation. To my knowledge no
        prior DLD-deformability study has imposed this kind of
        within-grid control.

  \item \textbf{Diagonal-subtracted angular separation
        $\Delta\theta_{\text{corr}}$:} I introduce a per-row baseline
        subtraction that removes the finite-$N$ noise floor from the
        off-diagonal cells, isolating the pure deformability-contrast
        signal. The diagonal cells deliver an empirically determined
        noise band of $\sim 1.3^\circ$ that is then quantitatively
        subtracted from every off-diagonal measurement.

  \item \textbf{Identification of a Capillary-number threshold:} the
        sweep reveals that stiffness-based sorting emerges only when
        the soft-species Capillary number exceeds
        $\mathrm{Ca}_{\text{soft}} \approx 0.2$. Below this
        threshold, the off-diagonal signal collapses to within the
        noise band regardless of stiffness ratio; above it,
        $|\Delta\theta_{\text{corr}}|$ grows monotonically up to a
        Capillary ratio of $\sim 13$ before saturating.

  \item \textbf{Practical chip-design rule:} the threshold
        $\mathrm{Ca}_{\text{soft}} > 0.2$ translates directly into a
        minimum flow speed for any DLD chip targeting a population
        of given stiffness. For physiological red blood cells in a
        $5\,\mu$m gap this corresponds to a centre-line velocity of
        order $1$~mm/s; for stiffer circulating tumour cells the
        required speed scales linearly with stiffness and may begin
        to conflict with shear-induced cell-damage limits --- a
        fundamental trade-off that this manuscript quantifies.

  \item \textbf{Provenance-rich open data:} every cell of the sweep
        ships a self-contained \verb|run_manifest.json| (git commit,
        compiler identity, resolved CXX flags, OpenMP thread count,
        full parameters, RNG seed) that is sufficient to reproduce
        the cell bit-for-bit when combined with the equations in
        the Methods section.
\end{itemize}

The manuscript directly addresses the scope of
\textit{Microfluidics and Nanofluidics}: it combines novel
numerical methodology, systematic parameter-space exploration, and
an immediately useful design rule for microfluidic cell-separation
chips. The approach is reproducible by any group with access to an
LBM-IBM solver of comparable fidelity, and the methodological
contributions (matched-membrane control, diagonal subtraction) are
transferable to any deformability-vs-geometry trade-off study in
related microfluidic separation techniques (inertial focusing,
viscoelastic migration, pinched-flow fractionation).

This is a single-author manuscript that I have prepared in full.
The work has not been published previously and is not under
consideration for publication elsewhere. I have no conflicts of
interest to declare. The work was conducted at the Department of
Natural Sciences, The Begum Nusrat Bhutto Women University,
Sukkur, with institutional support from the Sindh Higher Education
Commission.

Thank you for considering this manuscript. I look forward to your
response.

\closing{Sincerely,}

\end{letter}

\end{document}
